\scanpage{18}
\mspage{0}
% ll. 1–3
$\deleted{\unsure{P|\Gamma}}\Delta(E)={}$ groupe engendré par les éléments de la forme $[x,y]$ $(x,y)\in E$, avec les relations
\[
[x,x]=1,\quad [x,y][y,x]=1,\quad [x,y][y,z]=[x,z].
\]

% ll. 4–6
NB Je soupçonne que $\Delta(E)$ est isomorphe au sous-groupe du groupe libre $L(E)$ engendré par $E$, engendré par les éléments de la forme $[x,y]=x^{-1}y$.

% ll. 7–13
Si $\operatorname{card}E=n$ est fini, $\Delta(E)$ est un groupe libre à $n-1$ générateurs, si $E=[1,n]$, il est engendré par les générateurs libres $\deleted{\unsure{E_3|E_{\supset}|E_i}}$ $g_i=[i,i+1]$ pour $i\in[1,n-1]$. On a alors pour $i<j$
% chain at left, beside lines 8–10
\diagram{\begin{tikzcd}[ampersand replacement=\&,column sep=small]
1\arrow[r,dash]\&2\arrow[r,dash]\&3\arrow[r,dotted,dash]\&n-1\arrow[r,dash]\&n
\end{tikzcd}}
\[
\left\{\begin{aligned}
[i,j]&=g_ig_{i+1}\ldots g_{j-1}\\
[j,i]&=[i,j]^{-1}=g_{j-1}^{-1}\deleted{\unsure{g|j}}\ldots g_{i+1}^{-1}g_i^{-1}
\end{aligned}\right.
\]

% ll. 14–16
Résultat analogue si $E$ est dénombrable, alors $\Delta(E)$ est un groupe libre avec un ens dénombrable de générateurs libres.

% ll. 17–24
Soit $X$ une catégorie ayant un plus petit objet $e$, et un plus grand objet $a$. Alors $\pi_1(X,e)$ est isomorphe \deleted{\unsure{au|à un} \unsure{groupe|gr}} à un quotient du groupe $\Delta(\operatorname{Hom}(e,a))$, par les relations suivantes : pour tout système de doubles flèches
\diagram{\begin{tikzcd}[ampersand replacement=\&]
e\arrow[r,shift left=.6ex,"u"]\arrow[r,shift right=.6ex,"v"']\&x\arrow[r,shift left=.6ex,"\alpha"]\arrow[r,shift right=.6ex,"\beta"']\&a
\end{tikzcd}}
, \deleted{on} on a la relation
\[
\boxed{[\alpha v,\alpha u]=[\beta v,\beta u]}\quad\text{i.e.}\quad
[\alpha v,\alpha u][\beta v,\beta u]^{-1}=1.
\]
% archive sheet number: 17
